% ==============================================================================
% Extraits de Référence : Tableaux, Équations, Figures & Boîtes (mythesis.cls)
% ==============================================================================

% --- 1. Tableau professionnel (booktabs + tabularx + colonnes L/Y + siunitx) ---
\begin{table}[htbp]
  \centering
  \caption{Comparaison des performances expérimentales sur le jeu de test.}
  \label{tab:comparatif-modeles}
  \begin{tabularx}{\linewidth}{@{} l Y S[table-format=2.1] S[table-format=2.1] @{}}
    \toprule
    \textbf{Modèle} & \textbf{Architecture} & \textbf{Précision (\%)} & \textbf{Latence (\unit{\milli\second})} \\
    \midrule
    Référence A     & Transformeur (6 couches) & 84,2 & 12,4 \\
    Référence B     & Hybride CNN-RNN          & 81,7 &  9,1 \\
    \textbf{Proposé}& Bi-LSTM modulaire        & 89,5 &  7,8 \\
    \bottomrule
  \end{tabularx}
\end{table}

% --- 2. Équation alignée avec macros sémantiques (config/macros.sty) ---
\begin{align}
  \label{eq:objectif-optimisation}
  \mathcal{L}(\vect{\theta})
    &= \mathbb{E}_{x \sim \mathcal{D}}\!\left[-\log p_{\vect{\theta}}(x)\right] + \lambda\,\mathcal{R}(\vect{\theta}) \\
    &= \frac{1}{N}\sum_{i=1}^{N} \ell\bigl(f_{\vect{\theta}}(x_i), y_i\bigr) + \frac{\lambda}{2}\norm{\vect{\theta}}_2^2.
\end{align}

% --- 3. Figure vectorielle encadrée ---
\begin{figure}[htbp]
  \centering
  \begin{cadreFigure}[width=0.85\linewidth]
    \includegraphics[width=\linewidth]{figures/python-benchmark.pdf}
  \end{cadreFigure}
  \caption{Évolution comparée de la fonction de perte au fil des itérations.}
  \label{fig:benchmark-loss}
\end{figure}

% --- 4. Boîte de synthèse ---
\begin{boiteSynthese}{Point clé à retenir}
  Les macros \code{\textbackslash vect}, \code{\textbackslash norm} et les boîtes \code{boiteSynthese} s'adaptent automatiquement aux profils \code{screen} et \code{print}.
\end{boiteSynthese}
